\documentclass[12pt]{article}
\usepackage[utf8]{inputenc}
\usepackage{geometry}
\usepackage{graphicx}
\usepackage{hyperref}
\usepackage{listings}
\usepackage{xcolor}
\usepackage{enumitem}
\usepackage{media9} % For embedding videos

\geometry{margin=1in}

% Code listing style
\lstset{
    language=Java,
    basicstyle=\ttfamily\small,
    keywordstyle=\color{blue}\bfseries,
    commentstyle=\color{green!60!black},
    stringstyle=\color{red},
    numbers=left,
    numberstyle=\tiny\color{gray},
    stepnumber=1,
    numbersep=5pt,
    backgroundcolor=\color{gray!10},
    frame=single,
    breaklines=true,
    showstringspaces=false
}

\title{Chess Game - Phase 2 Documentation}
\author{Chess Game Project}
\date{\today}

\begin{document}

\maketitle

\section{Introduction}

This document describes the Phase 2 implementation of the Chess Game project, which includes a complete Graphical User Interface (GUI) built using Java Swing. The GUI provides an interactive chess playing experience with full chess rule validation, move history tracking, and customizable board settings.

\section{GUI Overview}

The chess game GUI consists of several key components:

\begin{itemize}
    \item \textbf{Chessboard Panel}: An 8x8 grid displaying the chessboard with alternating light and dark squares
    \item \textbf{Game History Panel}: Displays move history and captured pieces on the right side
    \item \textbf{Status Bar}: Shows whose turn it is and check/checkmate status
    \item \textbf{Menu Bar}: Provides access to game controls and settings
\end{itemize}

\subsection{Chessboard Display}

The chessboard is displayed as an 8x8 grid with alternating light and dark squares. Each chess piece is represented using Unicode chess symbols:
\begin{itemize}
    \item White pieces: \textcolor{white}{\textcolor{black}{♔ ♕ ♖ ♗ ♘ ♙}} (King, Queen, Rook, Bishop, Knight, Pawn)
    \item Black pieces: \textcolor{black}{♚ ♛ ♜ ♝ ♞ ♟} (King, Queen, Rook, Bishop, Knight, Pawn)
\end{itemize}

\subsection{Piece Movement}

The GUI supports two methods of moving pieces:

\begin{enumerate}
    \item \textbf{Click-and-Move}: Click on a piece to select it (highlighted in yellow), then click on the destination square
    \item \textbf{Drag-and-Drop}: Click and hold on a piece, drag it to the destination square, and release
\end{enumerate}

Both methods provide the same functionality and are fully validated according to chess rules.

\section{Chess Rules Implementation}

The GUI enforces all standard chess rules:

\subsection{Piece Movement Validation}

Each piece can only move according to its specific movement rules:
\begin{itemize}
    \item \textbf{Pawn}: Moves forward one square (two squares from starting position), captures diagonally
    \item \textbf{Rook}: Moves horizontally or vertically any number of squares
    \item \textbf{Knight}: Moves in an L-shape (two squares in one direction, one square perpendicular)
    \item \textbf{Bishop}: Moves diagonally any number of squares
    \item \textbf{Queen}: Combines rook and bishop movement (any direction)
    \item \textbf{King}: Moves one square in any direction
\end{itemize}

Invalid moves are rejected with error messages explaining why the move is not allowed.

\subsection{Check and Checkmate Detection}

The game detects and displays:
\begin{itemize}
    \item \textbf{Check}: When a king is under attack, the status bar displays "CHECK!"
    \item \textbf{Checkmate}: When a player has no legal moves to escape check, a pop-up dialog declares the winner
\end{itemize}

Players cannot make moves that would put their own king in check, ensuring proper chess gameplay.

\subsection{Capturing Pieces}

When a piece moves to a square occupied by an opponent's piece:
\begin{itemize}
    \item The opponent's piece is removed from the board
    \item The captured piece is displayed in the Game History Panel
    \item The move is recorded in the move history
\end{itemize}

\section{Extra GUI Features}

Three extra features have been implemented as required:

\subsection{Feature 1: Menu Bar with Game Controls}

The menu bar provides essential game management features:

\subsubsection{New Game}
\begin{itemize}
    \item Resets the board to the initial starting position
    \item Clears all move history and captured pieces
    \item Confirms with the user before resetting to prevent accidental data loss
\end{itemize}

\subsubsection{Save Game}
\begin{itemize}
    \item Allows players to save the current game state to a file
    \item Uses Java serialization to save:
    \begin{itemize}
        \item Complete board state (all piece positions)
        \item Current player's turn
        \item Full move history
    \end{itemize}
    \item Files are saved with a \texttt{.chess} extension
    \item Provides file chooser dialog for selecting save location
\end{itemize}

\subsubsection{Load Game}
\begin{itemize}
    \item Restores a previously saved game from a file
    \item Reconstructs the board, turn order, and move history
    \item Updates all GUI components to reflect the loaded game state
    \item Allows players to resume games exactly where they left off
\end{itemize}

\subsection{Feature 2: Settings Window for Customization}

The Settings window allows players to customize the chessboard appearance:

\subsubsection{Board Colors}
\begin{itemize}
    \item \textbf{Custom Colors}: Players can choose custom colors for light and dark squares using color pickers
    \item \textbf{Preset Themes}: Three preset color schemes are available:
    \begin{itemize}
        \item \textbf{Classic}: Traditional beige and brown colors (240, 217, 181 / 181, 136, 99)
        \item \textbf{Modern}: Green and cream colors (238, 238, 210 / 118, 150, 86)
        \item \textbf{Gray}: Modern gray theme (220, 220, 220 / 140, 140, 140)
    \end{itemize}
    \item Changes are applied immediately when "Apply" is clicked
\end{itemize}

\subsubsection{Board Size}
\begin{itemize}
    \item Adjustable board size using a slider (400-800 pixels)
    \item Real-time preview of selected size
    \item Automatically adjusts piece sizes to maintain proportions
\end{itemize}

\subsection{Feature 3: Game History Panel with Undo Button}

The Game History Panel provides comprehensive move tracking:

\subsubsection{Move History}
\begin{itemize}
    \item Displays all moves in chronological order
    \item Shows piece name, starting position, and destination position
    \item Uses standard algebraic notation (e.g., "E2 E4")
    \item Scrollable text area for long games
\end{itemize}

\subsubsection{Captured Pieces Display}
\begin{itemize}
    \item Shows captured pieces separately for White and Black
    \item Displays Unicode symbols of captured pieces
    \item Updates in real-time as pieces are captured
\end{itemize}

\subsubsection{Undo Functionality}
\begin{itemize}
    \item "Undo Move" button reverts the last move
    \item Restores the board to the previous state
    \item Restores captured pieces if applicable
    \item Updates move history and status display
    \item Maintains full undo history throughout the game
\end{itemize}

\section{Technical Implementation}

\subsection{Architecture}

The GUI is built using the Model-View-Controller (MVC) pattern:
\begin{itemize}
    \item \textbf{Model}: \texttt{Board} class manages game state
    \item \textbf{View}: GUI components (\texttt{ChessBoardPanel}, \texttt{GameHistoryPanel}) display the game
    \item \textbf{Controller}: \texttt{ChessGUI} handles user interactions and coordinates between components
\end{itemize}

\subsection{Key Classes}

\begin{description}
    \item[\texttt{ChessGUI}] Main window class that coordinates all GUI components
    \item[\texttt{ChessBoardPanel}] Custom JPanel that renders the chessboard and handles mouse interactions
    \item[\texttt{GameHistoryPanel}] Panel displaying move history and captured pieces
    \item[\texttt{SettingsDialog}] Dialog window for board customization
    \item[\texttt{GameState}] Serializable class for save/load functionality
    \item[\texttt{MoveRecord}] Records individual moves for history tracking
\end{description}

\subsection{Serialization}

For save/load functionality, the following classes implement \texttt{Serializable}:
\begin{itemize}
    \item \texttt{Board}
    \item \texttt{Piece} (and all subclasses)
    \item \texttt{Position}
    \item \texttt{GameState}
    \item \texttt{MoveRecord}
\end{itemize}

\section{Screenshots and Video Demonstration}

\subsection{Screenshots}

\textit{Note: Screenshots would be inserted here in the actual submission}

\subsubsection{Figure 1: GUI at Start of Game}
\textit{[Screenshot placeholder: Initial board setup with all pieces in starting positions]}

\subsubsection{Figure 2: Gameplay with Captured Pieces}
\textit{[Screenshot placeholder: Mid-game board with several moves made and pieces captured, showing captured pieces in history panel]}

\subsubsection{Figure 3: Checkmate and Winner Declaration}
\textit{[Screenshot placeholder: Checkmate position with pop-up dialog declaring the winner]}

\subsubsection{Figure 4: Settings Window}
\textit{[Screenshot placeholder: Settings dialog showing color customization and board size slider]}

\subsubsection{Figure 5: Game History Panel}
\textit{[Screenshot placeholder: Game History Panel showing move list and captured pieces]}

\subsection{Video Demonstration}

A screen recording demonstrating the chess game GUI in action is embedded below. The video shows:
\begin{itemize}
    \item Initial game setup and piece movement
    \item Click-and-move and drag-and-drop functionality
    \item Piece captures and move validation
    \item Check and checkmate detection
    \item Menu features (New Game, Save/Load)
    \item Settings customization
    \item Game History Panel and Undo functionality
\end{itemize}

% To embed a video, uncomment and modify the following:
% Place your video file (MP4 format recommended) in the same directory as this .tex file
% Then uncomment the lines below and replace "chess_game_demo.mp4" with your video filename

% Method 1: Using media9 package (works with Adobe Reader and some PDF viewers)
% \begin{center}
% \includemedia[
%   width=0.8\linewidth,
%   height=0.45\linewidth,
%   activate=pageopen,
%   addresource=chess_game_demo.mp4,
%   flashvars={source=chess_game_demo.mp4&autoPlay=false&loop=false}
% ]{\fbox{\parbox{0.8\linewidth}{\centering Click here to play video demonstration}}}{VPlayer.swf}
% \end{center}

% Method 2: Alternative - Link to video (if embedding doesn't work)
% \begin{center}
% \href{run:chess_game_demo.mp4}{\fbox{\parbox{0.8\linewidth}{\centering Click here to play video demonstration (requires external player)}}}
% \end{center}

\textit{[Video placeholder: Screen recording of chess game demonstration]}

\textbf{Note:} To embed your video:
\begin{enumerate}
    \item Record your screen using QuickTime (Mac), OBS, or similar software
    \item Export the video as MP4 format
    \item Place the video file in the same directory as this .tex file
    \item Uncomment the video embedding code above and replace \texttt{chess\_game\_demo.mp4} with your filename
    \item Compile the PDF - the video will be playable directly in the PDF viewer
\end{enumerate}

\section{Questions, Comments, and Explanations}

\subsection{Design Decisions}

\begin{itemize}
    \item \textbf{Unicode Symbols}: Chess pieces are displayed using Unicode characters for simplicity and cross-platform compatibility, rather than loading image files.
    
    \item \textbf{Move Validation}: All moves are validated before execution, including checking if the move would put the player's own king in check. This ensures proper chess gameplay.
    
    \item \textbf{Undo Implementation}: The undo feature maintains a history of board states rather than just reversing moves, ensuring accuracy even with complex game situations.
    
    \item \textbf{Save/Load Format}: Binary serialization was chosen over text-based formats for efficiency and to preserve the complete game state accurately.
\end{itemize}

\subsection{Challenges and Solutions}

\begin{itemize}
    \item \textbf{Check Detection}: Implementing check detection required careful consideration of piece movement rules and board state. The solution uses a test board copy to simulate moves before applying them.
    
    \item \textbf{Drag-and-Drop}: Implementing smooth drag-and-drop while maintaining piece validation required careful handling of mouse events and coordinate conversion.
    
    \item \textbf{Serialization}: Making all game classes serializable while maintaining backward compatibility required careful design of the class hierarchy.
\end{itemize}

\subsection{Future Enhancements}

Potential improvements for future versions:
\begin{itemize}
    \item Highlighting legal moves when a piece is selected
    \item Sound effects for moves and captures
    \item Timer functionality for timed games
    \item Network multiplayer support
    \item Opening book and move suggestions
    \item Export game to PGN (Portable Game Notation) format
\end{itemize}

\section{Conclusion}

The Phase 2 implementation successfully delivers a fully functional chess game GUI with:
\begin{itemize}
    \item Complete chess rule validation
    \item Intuitive user interface with multiple interaction methods
    \item Three comprehensive extra features
    \item Robust save/load functionality
    \item Professional appearance and user experience
\end{itemize}

The implementation follows object-oriented design principles and provides a solid foundation for future enhancements.

\end{document}

